

\spis{Ukryta moc defektów w~diamencie}
{Mariusz Mrózek}

\wtyt{Ukryta moc defektów w~diamencie}

\waut{Mariusz MRÓZEK*}
\aafil{Wydział Fizyki, Astronomii i~Informatyki Stosowanej, Uniwersytet Jagielloński}

Diamenty znane są jako najtwardsze naturalnie występujące substancje na Ziemi. Podziwiamy je za ich blask i~nieskazitelne piękno. Jednak to właśnie drobne defekty w~ich strukturze krystalicznej nadają im wyjątkowe właściwości, które czynią je niezwykle cennymi dla naukowców.

Idealny diament to taki, w~którym, każdy atom węgla jest ułożony w~perfekcyjnej sieci o~strukturze kubicznej. Wszelkie zmiany w~tej strukturze, takie jak wtrącenia innych pierwiastków, pęknięcia czy deformacje, nazywamy defektami. Choć na pierwszy rzut oka mogą wydawać się wadami, to w~rzeczywistości właśnie one stanowią o~wyjątkowych właściwościach diamentów.

Ciekawym przykładem wpływu defektów na właściwości tych kamieni jest ich przewodność elektryczna. Czysty diament jest izolatorem, co oznacza, że nie przewodzi prądu elektrycznego. Jednak obecność defektów, takich jak inkluzje atomów azotu (zastępujące część atomów węgla), mogą wprowadzić do struktury diamentu stany energetyczne, które umożliwiają przewodzenie prądu. To właśnie dzięki takim efektom diamenty są wykorzystywane w~elektronice, np.~w~tranzystorach~[1].\marg{{\includegraphics[clip,trim=33 23 20 20,width=4cm]{diament}}\\[4pt]
  Struktura diamentu. Źródło: Wikimedia Commons.}

Innym przykładem wpływu defektów na właściwości diamentów jest ich kolor. Idealny diament jest bezbarwny, ale obecność różnych defektów może nadać mu różne zabarwienie. Na przykład niebieski kolor diamentów jest wynikiem obecności atomów boru, a~czerwony wynikiem deformacji struktury kryształu. Te defekty kolorystyczne sprawiają, że diamenty są jeszcze bardziej poszukiwane i~cenne [2].

Najciekawszymi defektami punktowymi w~diamencie są centra barwne. To defekty struktury krystalicznej, które mają zdolność do absorpcji oraz emisji światła w~zakresie optycznym. Diamenty mają wysoką przezroczystość, co pozwala na łatwą obserwację fluorescencji lub absorpcji pochodzącej od różnych centrów barwnych, które w~nich się znajdują.
\marg{
Czym jest przerwa wzbroniona? Elektrony w~atomie mogą zajmować tylko określone stany energetyczne. Ze względu na zakaz Pauliego elektrony nie mogą znajdować się w~tym samym stanie. W~konsekwencji w~atomach elektrony zapełniają różne powłoki energetyczne (orbitale), zaczynając od najniższych energii. W~każdym stanie znajdują się co najwyżej dwa elektrony, różniące się spinem. Kryształ składa się z~ogromnej liczby atomów. Z~powodu ich ciasnego upakowania występuje częściowe przekrywanie się orbitali i~przesuwanie odpowiadających im energii. Powstaje wtedy tzw. \emph{pasmo energetyczne}. W~krysztale odpowiednikiem jednego stanu energetycznego w~atomie jest całe pasmo złożone z~ogromnej ilości stanów o~zbliżonych do siebie energiach. W~różnych kryształach przerwy pomiędzy pasmami są większe lub mniejsze, a~czasami przerwy może nie być -- pasma mogą się częściowo pokrywać. W~czystym diamencie przerwa (tzw. \emph{przerwa wzbroniona}) między ostatnim zapełnionym pasmem a~pierwszym niezapełnionym jest duża. Właśnie dlatego czysty diament jest izolatorem -- nie ma w~nim swobodnych nośników, które mogłyby przewodzić prąd. Przewodzeniem prądu zajmują się elektrony z~pasma częściowo tylko zapełnionego -- jeżeli takiego pasma nie ma, to mamy izolator. Więcej o~pasmach i~przewodnictwie można przeczytać w~$\Delta^{7}_{22}$, w~artykule \href{https://www.deltami.edu.pl/media/articles/2022/07/delta-2022-07-dioda-bez-dziur.pdf}{\emph{Dioda bez dziur}}.
}
Diament ma szeroką przerwę wzbronioną, natomiast poziomy energetyczne elektronów związanych z~takim centrum barwnym są dobrze oddalone od pasm diamentu (znajdują się w~przerwie między pasmami), co pozwala traktować centra barwne jako ,,nowe atomy'' uwięzione w~sieci krystalicznej diamentu i~dobrze odizolowane od wpływu otoczenia. Dotychczas w~diamencie zostało zaobserwowanych ponad 600 różnych centrów barwnych.
Centrum barwne azot-wakancja $NV^{-}$ to jeden z~wielu defektów sieci krystalicznej diamentu. Składa się on z~azotu~\textit{N} (zastępuje węgiel), w~którego najbliższym sąsiedztwie znajduje się wakancja~\textit{V}, czyli węzeł sieci nieobsadzony atomem (w~miejscu, gdzie powinien być atom, jest pusta przestrzeń). Atom azotu $^{14}N$ posiada pięć elektronów walencyjnych. Trzy z~nich są związane z~atomami węgla znajdującymi się wokół niego. Pozostałe dwa skierowane są do wakancji. Układ taki nazywany jest neutralnym centrum barwnym $NV^0$. Neutralnie naładowane centrum $NV^0$ wydaje się obiecującym kandydatem do zastosowań w~systemach informacji kwantowej, ze względu na posiadanie spinu $(S = 1/2)$. Jeśli do wakancji przyłączony zostanie kolejny elektron, to defekt ten zmienia się w~centrum barwne $NV^{-}$, któremu więcej uwagi poświęcimy poniżej. Rzadziej występuje sytuacja, gdy są tylko cztery
elektrony, wówczas mamy do czynienia z~centrum barwnym dodatnio naładowanym $NV^+$ [3].

Dla centrów barwnych $NV^{-}$ stosowana jest podstawowa metoda badawcza nazwana optycznie wykrywanym rezonansem magnetycznym (ODMR, Optically Detected Magnetic Resonance). Polega na obserwacji zmiany natężenia czerwonej fluorescencji próbki wzbudzanej
zielonym laserem podczas skanowania częstości pola mikrofalowego wokół rezonansu
magnetycznego (około 2,87~GHz). ODMR umożliwia wykonywanie pomiarów wielu wielkości fizycznych, jak pole magnetyczne i~elektryczne, temperaturę czy ciśnienie.

\baselineskip12pt minus.1pt
\marg[-1]{{\bf \hypertarget{odp}{}\hyperlink{pytania}{Odpowiedzi na pytania ze s.~12}.}
\begin{enumerate}[label=(\alph*)]
\item Między 62 ($\mathtt{sixty\ two}$) i~13~($\mathtt{thirteen}$).
\item Pierwsza to 40 ($\mathtt{czterdzieści}$), a~przedostatnia to 13 ($\mathtt{trzynaście}$).
\end{enumerate}}Wartości pola magnetycznego mierzone z~wykorzystaniem $NV^{-}$ z~czułością sięgającą nawet setek $pT/\sqrt{Hz}$ (jednostka porównawcza dla wielu technik pomiaru wartości pola magnetycznego -- im mniejsza wartość, tym bardziej czułe urządzenie). Taka dokładność pozwala na rejestrację np. sygnałów magnetycznych pochodzących od serca na powierzchni klatki piersiowej [4].
Warto w~tym miejscu wspomnieć, że dla centrum  barwnego azot-wakancja nie wszystkie pomiary indukcji pola magnetycznego wymagają wykorzystywania techniki ODMR.  Wraz ze wzrostem przyłożonego pola magnetycznego, aż do około 30~mT, natężenie fluorescencji centrów $NV^{-}$  spada o~kilkanaście procent. Po odpowiednim wykalibrowaniu zależność ta może być wykorzystana do wyznaczania wartości pola magnetycznego, oferując szeroki zakres dynamiczny pola.
Dokładne pomiary spektroskopowe pokazały, że położenie rezonansu ODMR, pochodzącego od centrum barwnego $NV^{-}$, w~diamencie zależy także od temperatury otoczenia. Zależność ta jest wynikiem rozszerzalności cieplnej oraz oddziaływania elektronów z~fononami. Rezonans ODMR przesuwa się ku niższym częstotliwościom wraz ze wzrostem temperatury.  Dzięki temu zjawisku centrum barwne $NV^{-}$ może pełnić rolę miniaturowego termometru. Dla niedużego zakresu zmian temperatur zmiana częstotliwości rezonansu może być w~przybliżeniu traktowana jako liniowa, a~w~temperaturach bliskich pokojowej przesunięcie to wynosi około $-74$~kHz/K.

Przedstawione do tej pory zastosowania czujników z~centrami $NV^{-}$ obejmują także nanodiamenty. Możliwości wyznaczania pól i~temperatur w~skali nano są ogromnym atutem umożliwiającym pomiary tych  wielkości fizycznych z~wysoką przestrzenną zdolnością rozdzielczą. W~połączeniu  z~nowoczesnymi technikami współczesnej mikroskopii optycznej i~mikroskopii skaningowej otwiera to bardzo atrakcyjne perspektywy dla wielu dyscyplin nauki. Nanodiament zawierający centrum barwne $NV^{-}$ składa się z~nietoksycznych związków: węgla – podstawowego budulca związków organicznych i~niewielkich ilości azotu, idealnie nadaje się więc do badań biologicznych. Taki nanodiament może być umieszczony w~komórce biologicznej i~służyć jako marker biologiczny lub czujnik temperatury. Wykorzystuje się także fakt, że centra barwne $NV^{-}$ nie ulegają blaknięciu pod wpływem światła (\textit{photo bleaching}). Dzięki temu diamenty z~centrami $NV^{-}$ nadają się do długoterminowego fluorescencyjnego znacznikowania i~obrazowania przyłączonych do nich związków chemicznych czy biologicznych [5].

\marg[36]{\setbox0=\hbox{[0]\enskip}\def\pp[#1]{\vskip2pt\hangindent\wd0\noindent\rlap{[#1]}\kern\wd0\ignorespaces}
  Literatura
  \pp
  [1]Araujo D., Suzuki M., Lloret F., Alba G., Villar P., ``Diamond for Electronics: Materials, Processing and Devices'', Materials, 14(22), (2021).
  \pp
  [2] Rachminov E., ``The Fancy Color Diamond Book: Facts and Secrets of Trading in Rarities'', New York: Diamond Odyssey (2009).
  \pp
  [3] Jensen K., Kehayias P., Budker D., ``Magnetometry with Nitrogen-Vacancy Centers in Diamond'', Smart Sensors, Measurement and Instrumentation, 19, Springer (2017). \pp
  [4] Arai K., Kuwahata A., Nishitani D. et~al., ``Millimetre-scale magnetocardiography of living rats with thoracotomy'', Commun. Phys. 5, 200 (2022).
  \pp
  [5] Schirhagl R, Chang K, Loretz M, Degen CL., ``Nitrogen-vacancy centers in diamond: nanoscale sensors for physics and biology'', Annu. Rev. Phys. Chem. (2014).
  \pp
  [6] Gulka M., Wirtitsch D., Ivády V. et~al., ``Room-temperature control and electrical readout of individual nitrogen-vacancy nuclear spins'', Nat. Commun. 12, 4421 (2021).}
\looseness-1{\spaceskip.33em minus.05emZ~naukowego punktu widzenia diament z~centrami barwnymi azot-wakancja wydaje się idealnym materiałem do zbudowania urządzenia wykorzystującego właściwości mechaniki kwantowej. W~świecie kwantów diament naprawdę się wyróżnia, a~jego właściwości optyczne doskonale nadają się do implementacji w~nowych technologiach kwantowych. Centra barwne azot-wakancja w~diamencie oferują możliwość budowania indywidualnych systemów kwantowych, które działają dobrze w~temperaturze pokojowej. Operacje na qubitach, czyli superpozycjach stanów spinowych $m_S$, można realizować za pomocą impulsów mikrofalowych w~centrach barwnych $NV^{-}$. Możliwość indywidualnego adresowania wybranego miejsca próbki, a~nawet pojedynczych centrów, za pomocą konwencjonalnej mikroskopii i~zastosowania  impulsów mikrofalowych pozwalają na pełną kwantową kontrolę stanu centrum barwnego.  Dla dobrego wykorzystania qubitu, pozwalającego na wykonanie odpowiednio dużej liczby operacji, konieczne jest z jednej strony osiąganie dostatecznej szybkości zmiany stanu (wysokiej częstości Rabiego), a~z~drugiej strony dostatecznie długiego jego czasu życia (wolnej relaksacji).
  % Dla dobrego wykorzystania qubitu, pozwalającego na wykonanie odpowiednio dużej liczby operacji, konieczne jest z~jednej strony osiąganie dostatecznej szybkości zmiany stanu (wysokiej częstości Rabiego) i~dostatecznie długiego jego czasu życia (wolnej relaksacji), z~drugiej strony.
  Centra barwne $NV^{-}$ w~diamencie zapewniają spełnienie obu tych warunków. Co ważne, nie jest konieczne stosowanie skomplikowanych technik kriogenicznych i~mogą one być wykorzystywane nawet w~temperaturze pokojowej, co jest wielkim atutem. Trudno takie efekty osiągnąć z~innymi kandydatami na qubity z~zakresu fizyki ciała stałego~[6].}


Podsumowując, defekty w~diamencie, choć na pierwszy rzut oka mogą wydawać się wadami, w~rzeczywistości nadają im wyjątkowe właściwości, które czynią je niezwykle cennymi. Bez tych drobnych niedoskonałości diamenty nie byłyby dziś tak fascynującymi i~wartościowymi kamieniami szlachetnymi.




\normalsize